% ============================================================
%  Progressão Geométrica — Apostila Premium | PratiqueJá
%  Engine: XeLaTeX
% ============================================================
\documentclass[12pt]{article}
\usepackage{../../pratiqueja}

\newcommand{\hlm}[1]{{\color{darkblue}#1}}

\setsubject{Progressão Geométrica}

\begin{document}
\pagenumbering{arabic}
\setstretch{1.2}
\color{bodytext}

\heroheader{Progressão Geométrica (PG)}%
           {Termo geral, soma finita, PG infinita e interpolação}%
           {Matemática · Avançado}

\setlength{\columnsep}{18pt}
\setlength{\columnseprule}{0.5pt}
\def\columnseprulecolor{\color{exborder}}

\begin{multicols}{2}

\TW{Def}{darkblue}{Definição}{O que é uma Progressão Geométrica?}

\regra{
Uma \hl{Progressão Geométrica (PG)} é uma sequência em que cada termo, a partir do segundo, é obtido multiplicando o anterior por uma constante chamada \hl{razão $q$}.

\[
  q = \frac{a_n}{a_{n-1}}, \quad \forall\; n > 1
\]
}

\exemp{
\textbf{Crescente} ($q > 1$): $2, 6, 18, 54, \ldots$ com $q = 3$

\textbf{Decrescente} ($0 < q < 1$): $16, 8, 4, 2, \ldots$ com $q = \frac{1}{2}$

\textbf{Alternante} ($q < 0$): $2, -6, 18, -54, \ldots$ com $q = -3$

\textbf{Constante} ($q = 1$): $5, 5, 5, 5, \ldots$
}

\TW{$a_n$}{darkblue}{Termo Geral}{Calculando qualquer termo sem listar todos}

\regra{
\[
  a_n = a_1 \cdot q^{n-1}
\]
}

\exemp{
\textbf{PG $3, 6, 12, 24, \ldots$ — qual o 10º termo?}

$a_1 = 3$, $q = 2$

$a_{10} = 3 \cdot 2^9 = 3 \cdot 512 = \hlm{1536}$

\textbf{$a_1 = 5$, $q = 2$. Para qual $n$ temos $a_n = 160$?}

$160 = 5 \cdot 2^{n-1} \Rightarrow 2^{n-1} = 32 = 2^5$

$\hlm{n = 6}$ — o termo 160 é o 6º.

\textbf{$a_2 = 6$ e $a_5 = 48$. Determine $a_1$ e $q$:}

$a_5 = a_2 \cdot q^3 \Rightarrow 48 = 6q^3 \Rightarrow q^3 = 8 \Rightarrow q = 2$

$a_1 = a_2 / q = 6/2 = \hlm{3}$

PG: $3, 6, 12, 24, 48, \ldots$
}

\TW{$S_n$}{darkblue}{Soma dos Termos}{Soma dos $n$ primeiros termos da PG}

\regra{
Para $q \neq 1$:
\[
  S_n = \frac{a_1 \cdot (q^n - 1)}{q - 1}
\]
Quando $a_n$ é conhecido: $\;\; S_n = \dfrac{a_n \cdot q - a_1}{q - 1}$
}

\nota{\textbf{Caso $q = 1$:} todos os termos são iguais a $a_1$, portanto $S_n = a_1 \cdot n$.}

\exemp{
\textbf{Soma dos 6 primeiros termos da PG $2, 6, 18, 54, \ldots$:}

$a_1 = 2$, $q = 3$

$S_6 = \dfrac{2 \cdot (3^6 - 1)}{3 - 1} = \dfrac{2 \cdot 728}{2} = \hlm{728}$
}

\TW{$S_\infty$}{badgepurple}{PG Infinita Decrescente}{Soma de infinitos termos quando $|q| < 1$}

\regra{
Quando $|q| < 1$, os termos se aproximam de zero e a soma converge:
\[
  S_\infty = \frac{a_1}{1 - q}
\]
}

\nota{\textbf{Condição:} a fórmula $S_\infty$ só vale quando $|q| < 1$. Para $|q| \geq 1$, a soma diverge.}

\exemp{
\textbf{PG $8, 4, 2, 1, \frac{1}{2}, \ldots$:}

$a_1 = 8$, $q = \frac{1}{2}$

$S_\infty = \dfrac{8}{1 - \frac{1}{2}} = \dfrac{8}{\frac{1}{2}} = \hlm{16}$

\textbf{Dízima $0{,}\overline{3}$ como fração:}

$0{,}\overline{3} = 0{,}3 + 0{,}03 + \ldots$ \;\; $a_1 = 0{,}3$, $q = 0{,}1$

$S_\infty = \dfrac{0{,}3}{1 - 0{,}1} = \dfrac{0{,}3}{0{,}9} = \hlm{\dfrac{1}{3}}$
}

\TW{Int}{badgepurple}{Interpolação Geométrica}{Inserir meios geométricos entre dois termos}

\regra{
Interpolar $k$ \hl{meios geométricos} entre $a$ e $b$ (com $a, b > 0$) cria uma PG com $k+2$ termos. A razão é:
\[
  q = \sqrt[k+1]{\frac{b}{a}} = \left(\frac{b}{a}\right)^{\!\frac{1}{k+1}}
\]
}

\exemp{
\textbf{2 meios entre $3$ e $24$:}

$q = \sqrt[3]{\frac{24}{3}} = \sqrt[3]{8} = 2$

$\hlm{\text{PG: } 3,\; 6,\; 12,\; 24}$

\textbf{3 meios entre $1$ e $16$:}

$q = \sqrt[4]{16} = 2$

$\hlm{\text{PG: } 1,\; 2,\; 4,\; 8,\; 16}$ \quad (5 termos)
}

\nota{\textbf{Atenção:} para $k$ par, $a$ e $b$ devem ter o mesmo sinal para que $q$ seja real.}

\TW{Prop}{badgepurple}{Propriedades}{Relações entre os termos de uma PG}

\exemp{
\textbf{Média geométrica:} qualquer termo é a média geométrica dos vizinhos:
\[
  a_k^2 = a_{k-1} \cdot a_{k+1}
\]
Ex: $2, \mathbf{6}, 18$: $\quad 6^2 = 36 = 2 \times 18$ \;\ok

\textbf{Produto equidistante:} constante entre extremos:
\[
  a_k \cdot a_{n-k+1} = a_1 \cdot a_n
\]
Ex: $1, 2, 4, 8, 16$: $\;\;a_2 \cdot a_4 = 16 = a_1 \cdot a_5$ \;\ok

\textbf{Logaritmos:} se $(a_1, a_2, \ldots)$ é PG, então $(\log a_1, \log a_2, \ldots)$ é PA com razão $\log q$.

\textbf{Forma alternativa:} $a_n = a_k \cdot q^{n-k}$ (útil quando $a_k$ é mais conveniente que $a_1$)
}

\TW{Prob}{badgepurple}{Problemas Contextualizados}{Aplicação da PG em situações reais}

\exemp{
\textbf{Prob. 1 — Juros compostos:}

R\$\,1.000,00 a 10\% a.m. por 5 meses.

$a_1 = 1000$, $q = 1{,}10$, $n = 6$ (mês 0 ao mês 5)

$M = 1000 \cdot (1{,}10)^5 \approx \hlm{\text{R\$}\;1.610{,}51}$

\textbf{Prob. 2 — Bactérias:}

500 bactérias, $q = 2$ por hora.

$a_9 = 500 \cdot 2^8 = 500 \cdot 256 = \hlm{128.000}$ bactérias após 8 h.

\textbf{Prob. 3 — Dízima periódica:}

$0{,}\overline{27} = 0{,}27 + 0{,}0027 + \ldots$

$a_1 = 0{,}27$, $q = 0{,}01$

$S_\infty = \dfrac{0{,}27}{0{,}99} = \dfrac{27}{99} = \hlm{\dfrac{3}{11}}$
}

\end{multicols}

\end{document}
